% Archivo generado por bd/generar_tablas.ps1 a partir de bd/fase_posterior/crear_tablas_fase_posterior.sql. No editar a mano.
\subsubsection*{
Relaciones entre jugadores y tutorial
}
\addcontentsline{toc}{subsubsection}{
Relaciones entre jugadores y tutorial
}

\atributosde{Amistad}{Amistad recíproca entre dos cuentas, guardada una sola vez por par ordenado (\mbox{CU-31} \mbox{RN-03}).}
\begin{tablaatributos}
\texttt{id\_\allowbreak{}usuario\_\allowbreak{}a} & \texttt{INT} & No & PK, FK & El menor de los dos identificadores. \\ \hline
\texttt{id\_\allowbreak{}usuario\_\allowbreak{}b} & \texttt{INT} & No & PK, FK & El mayor de los dos identificadores. \\ \hline
\texttt{fecha\_\allowbreak{}amistad} & \texttt{DATETIME2(0)} & No &  & Momento en que se aceptó la solicitud. Por omisión, la fecha actual. \\ \hline
\end{tablaatributos}

\atributosde{Solicitud}{Solicitud de amistad pendiente. Tiene dirección, a diferencia de la amistad, y se borra al responderse (\mbox{CU-30}, \mbox{CU-31}).}
\begin{tablaatributos}
\texttt{id\_\allowbreak{}solicitud} & \texttt{INT}, identidad & No & PK & Identificador de la solicitud. \\ \hline
\texttt{id\_\allowbreak{}solicitante} & \texttt{INT} & No & FK & Jugador que la envía. \\ \hline
\texttt{id\_\allowbreak{}destinatario} & \texttt{INT} & No & FK & Jugador que la recibe. \\ \hline
\texttt{fecha\_\allowbreak{}solicitud} & \texttt{DATETIME2(0)} & No &  & Envío; caduca a los treinta días (\mbox{CU-31} \mbox{RN-05}). Por omisión, la fecha actual. \\ \hline
\end{tablaatributos}

\atributosde{ProgresoTutorial}{Avance de cada cuenta en cada lección del tutorial; es la relación N:M entre Usuario y Leccion (\mbox{CU-41}).}
\begin{tablaatributos}
\texttt{id\_\allowbreak{}usuario} & \texttt{INT} & No & PK, FK & Cuenta. \\ \hline
\texttt{id\_\allowbreak{}leccion} & \texttt{TINYINT} & No & PK, FK & Lección. \\ \hline
\texttt{paso\_\allowbreak{}actual} & \texttt{TINYINT} & No &  & Último paso completado. \\ \hline
\texttt{fecha\_\allowbreak{}actualizacion} & \texttt{DATETIME2(0)} & No &  & Último avance. Por omisión, la fecha actual. \\ \hline
\texttt{fecha\_\allowbreak{}completada} & \texttt{DATETIME2(0)} & Sí &  & Momento en que se terminó la lección. \\ \hline
\texttt{completada} & calculada & --- &  & Si la lección está terminada. Calculada: es exactamente que \texttt{fecha\_\allowbreak{}completada} no sea nula. \\ \hline
\end{tablaatributos}

